%Why I study physics
How do simple local interactions give rise to collective behavior that no single constituent possesses? I first met this question as \textit{emergent phenomena} (``More is Different''), and it has shaped my taste for doing research and the topic I have been working on since undergraduate. I want to pursue a PhD in physics because I want to work on such questions independently, in condensed matter theory, where ideas from quantum information, gauge theory, and topology meet. My long-term goal is a career in academic research, where I can keep working on quantum matter and contributing to this field.

%How I got started
I found my field when I attended Prof.~Yi-Ping Huang's talk about quantum spin liquids, phases of matter that have no local order parameter and are defined by long-range entanglement and fractionalized excitations. I joined his group as an undergraduate. My first project was to express the ground state of the Kitaev honeycomb model as a tensor network, a representation that is numerically efficient and also exposes a state's entanglement structure. I had to teach myself frustrated magnetism and tensor networks to do this. The topic did not come to me finished. It took shape through months of discussion with my advisor, and I learned that choosing a good question and explaining why it matters are skills of their own. The project earned me an undergraduate research scholarship from Taiwan's National Science and Technology Council (NSTC).

%Research accomplishments
I stayed in the group for my master's degree and began collaborating with Dr.~Sungkit Yip at Academia Sinica, whom Prof.~Huang introduced to me during a 2024 summer internship. We first studied how topological order in the Kitaev ground state arises from many-body entanglement. Dr.~Yip's rule of \textit{always start from a simple example} guided this work. With this spirit, we started from a two-spin model, moved to a one-dimensional chain, and then extended the picture to the full honeycomb lattice. I also used symmetry-resolved exact diagonalization and the Lanczos method, which I learned in my courses, to test my analytical results numerically.

We then moved to the density matrix itself, and this became the main result of my master's work, a paper on which I am first author, published in \textit{Physical Review B}. The standard calculation of Kitaev-model entanglement is done in an enlarged Majorana Hilbert space with gauge redundancy, which hides the physical origin of the results. We worked directly in the physical spin Hilbert space and derived the exact density matrix of the Kitaev spin liquid in terms of spin-string operators, organized by equivalence classes that depend only on the strings' endpoints. Combining this structure with the Gauss law of the emergent $\mathbb{Z}_2$ gauge theory and the model's exact 1-form Wilson symmetry, we showed that the reduced density matrix splits into $2^{|\partial A|/2-1}$ identical blocks. This solves the entanglement spectrum exactly, and it explains both its extensive degeneracy and the gauge part of the entanglement entropy. [State your own contribution in two or three sentences: which derivations, proofs, numerics, or key ideas were yours.] To do this I had to learn lattice gauge theory, higher-form symmetries, and quantum-information tools for spin systems, and apply them to a real problem.

%Challenges and setbacks
Research rarely moves in a straight line, and this project showed me that. Our first direction, the entanglement structure of the ground state, did not yet answer what we wanted to know, so we changed the question and turned to the density matrix, even though it meant setting aside work already done. A second obstacle came later. The existing Majorana method for computing the entropy seemed to apply only when a subsystem has an even number of sites, since an odd number of Majorana fermions does not define a Hilbert space. It took us a long time to see that the problem was in the Majorana description and not in the physics. Working in the spin representation, we found that the odd case involves a subsystem fermion parity, which gives an extra twofold degeneracy in the entanglement spectrum, and we could describe it with an explicit dimer-covering operator. Some problems also remain open. We did not find a 't~Hooft symmetry operator that commutes with the Hamiltonian and is independent of the Wilson loops, and we said so in the paper. Setbacks like these taught me to pursue several approaches at once, to check each step against a simple limit, and to state clearly what has and has not been shown. {\color{red} real personal or academic setback?}

%Current interests and next steps 
My current work builds on this framework. I am studying quantum magic, measured by the stabilizer R\'enyi entropy, in the Kitaev model, starting with the one-dimensional chain, and I have been building a numerical library for it. The density-matrix and string-operator approach suggests a way to study the magic of reduced density matrices. More broadly, I want to learn which structures of quantum spin liquids survive at finite temperature or in mixed states, where the entanglement entropy no longer plays its familiar role and the density matrix itself is the natural object. In graduate school I want to widen my toolkit beyond exactly solvable models to numerical methods for strongly interacting systems and to other topological and correlated phases.

{\color{red} Why CU Boulder}
%[Add 3--5 sentences. Name two or three faculty whose work connects to yours, with one specific topic or paper each, and mention any department strengths or collaborations that matter to you. Check the current faculty and research-group pages before writing this.]

In summary, I come to graduate study with research experience from the first idea to a published paper, with a record of finishing hard theoretical work, and with a clear picture of what I want to learn next. I want to continue that work as a PhD student and, eventually, to lead my own research program.